% ABSTRACT - Auto-updated section
% Last generated: 2026-01-15T05:34:34.872Z
% Source: Ollama deepseek-r1:32b

Decentralized Autonomous Organizations (DAOs) represent a paradigm shift in organizational governance, yet their design space remains poorly understood. We present a comprehensive multi-agent simulation framework for analyzing governance mechanisms in DAOs, enabling systematic study of how design parameters affect collective decision-making outcomes. Our framework models heterogeneous agent populations with varying participation strategies, token distributions, and delegation behaviors, supporting multiple voting mechanisms including token-weighted, quadratic, and conviction voting.

Through Monte Carlo simulations (21402 runs across 10 experiment sets), we investigate participation and capture dynamics in DAO governance: how quorum thresholds interact with turnout, how participation scales with DAO size, and how governance capture mitigation trades off against throughput.

We test hypotheses about quorum sensitivity, participation decay, and capture mitigation, and report results where experiments have completed. The paper is a living document; claims are updated as additional runs are incorporated. We release our simulation framework as open-source software to enable reproducible governance research and evidence-based mechanism design.

% === AUTO-GENERATION METADATA ===
% This section can be regenerated by running:
%   npm run paper:update -- --section abstract
% Using results from: {{RESULTS_DIR}}
